\documentclass[11pt]{article}
\usepackage[a4paper,margin=25mm]{geometry}
\usepackage[T1]{fontenc}
\usepackage{lmodern,microtype,amsmath,amssymb,amsthm,mathtools}
\usepackage[hidelinks]{hyperref}
\setlength{\emergencystretch}{2em}
\newtheorem{theorem}{Theorem}
\newtheorem{lemma}{Lemma}
\newcommand{\R}{\mathbb R}
\newcommand{\pos}{\succeq0}
\title{A second-order cone representation of the angular determinant}
\author{}
\date{23 October 2026}
\begin{document}
\maketitle
\vspace{-2em}

The angular determinant of a canonical Bell inequality contains eight
nonconstant Laurent monomials. Their coefficients are nonnegative;
only the constant can be negative. Positivity is therefore determined by
how much constant each balanced group of exponents can support. There are
twelve such elementary groups. We give their complete matrix representation,
including vanishing coefficients and minima approached at infinity.

Put $E=\{-1,0,1\}^2\setminus\{0\}$ and write $z^e=U^{e_1}V^{e_2}$.
The cone to be represented is
\begin{equation}\label{eq:cone}
 \mathcal L=\{(h,q):q_e\ge0,\quad
             h+\sum_{e\in E}q_ez^e\ge0\quad(U,V>0)\}.
\end{equation}
Its nine coordinates are $h$ and the eight $q_e$.

\section{The twelve circuits and their matrices}

Opposite pairs $\{e,-e\}$ give four circuits. Four more are
$\{(s,t),(-s,0),(0,-t)\}$, with $s,t\in\{-1,1\}$.
Their exponents add to zero. The remaining four have the form
$\{e,f,g\}$, where $e,f$ are adjacent corners of the square and
$e+f+2g=0$. Denote the eight triangles by $\mathcal T$ and all
twelve circuits by $\mathcal C$.

For each triangle $C$ allocate three coefficients $r_{C,e}$ to its
three monomials. The opposite pairs receive everything left over:
\begin{equation}\label{eq:residual}
 R_e=q_e-\sum_{C\in\mathcal T:\,e\in C}r_{C,e}.
\end{equation}
Each circuit also receives a constant $s_C$. The matrices below ensure
that its allocated monomials dominate $s_C$ for every positive $U,V$.

For a pair $C=\{e,-e\}$ require
\begin{equation}\label{eq:pair}
 \begin{pmatrix}2R_e&s_C\\s_C&2R_{-e}\end{pmatrix}\pos.
\end{equation}
For a triangle whose exponents add to zero, abbreviate its three
allocations by $r_1,r_2,r_3$, put $t=s_C/3$, and require
\begin{equation}\label{eq:equal}
 \begin{pmatrix}r_1&x\\x&r_2\end{pmatrix}\pos,\qquad
 \begin{pmatrix}r_3&y\\y&t\end{pmatrix}\pos,\qquad
 \begin{pmatrix}x&t\\t&y\end{pmatrix}\pos.
\end{equation}
For $e+f+2g=0$, use $r_1,r_2,r_3$ in that order, put $t=s_C/4$,
and require
\begin{equation}\label{eq:weighted}
 \begin{pmatrix}r_1&x\\x&r_2\end{pmatrix}\pos,\qquad
 \begin{pmatrix}x&t\\t&r_3/2\end{pmatrix}\pos.
\end{equation}
Finally impose the single scalar inequality
\begin{equation}\label{eq:constant}
                         h+\sum_{C\in\mathcal C}s_C\ge0.
\end{equation}
Every matrix entry is affine in the nine inputs and the auxiliary
variables. The allocations and $R_e$ are automatically nonnegative
because they occur on matrix diagonals.

\begin{theorem}\label{thm:lift}
Equations~\eqref{eq:residual}--\eqref{eq:constant} represent
$\mathcal L$ exactly. Their block diagonal pencil has 24 blocks of
order two and one scalar block, hence total order 49. It uses 48
auxiliary variables and rational coefficients.
\end{theorem}

\begin{proof}
We first show that the matrices imply positivity. For an opposite
pair, \eqref{eq:pair} gives
$s_C\le2\sqrt{R_eR_{-e}}$. The arithmetic--geometric mean inequality
then gives $R_ez^e+R_{-e}z^{-e}\ge s_C$.

In \eqref{eq:equal}, all diagonal entries are nonnegative and
$x^2\le r_1r_2$, $y^2\le r_3t$, $t^2\le xy$.
If $t>0$, these imply $t^3\le r_1r_2r_3$; if $t=0$ the desired
bound is immediate. Since the three monomials have product one,
their allocated sum is at least
$3(r_1r_2r_3)^{1/3}\ge s_C$.

Similarly, \eqref{eq:weighted} implies
$t^4\le r_1r_2(r_3/2)^2$. Applying the arithmetic--geometric mean
inequality to $r_1z^e,r_2z^f,(r_3/2)z^g,(r_3/2)z^g$ shows that
their sum is at least $4|t|\ge s_C$. Thus negative $s_C$ causes no
problem in either a pair or a weighted triangle; the ordinary triangle
already forces $s_C\ge0$.
Summing all twelve bounds and using \eqref{eq:residual} and
\eqref{eq:constant} proves membership in \eqref{eq:cone}.

For the converse, suppose first that $h\ge0$. Set all triangle
allocations, constants and auxiliary variables to zero. The pair
blocks are nonnegative diagonal matrices, and \eqref{eq:constant}
holds. Now let $h<0$ and let $I=\{e:q_e>0\}$. Positivity forces
$0\in\operatorname{conv} I$: otherwise a separating direction in
$(\log U,\log V)$ makes every active monomial tend to zero.

Let $J$ be the active exponents on the smallest face of
$\operatorname{conv}I$ containing zero. An exposing direction sends
all off-face monomials to zero and leaves the face monomials fixed.
Hence $h+\sum_{e\in J}q_ez^e$ is nonnegative. In the span of $J$,
the exponential sum is strictly convex and tends to infinity along
every unbounded ray. It has a positive minimizer $z_c=(U_c,V_c)$,
unique modulo directions orthogonal to $J$. Put
\[
 w_e=q_ez_c^e\quad(e\in J),\qquad w_e=0\quad(e\notin J).
\]
Stationarity and positivity at this minimum give
\begin{equation}\label{eq:balance}
               \sum_e w_e e=0,\qquad h+\sum_e w_e\ge0.
\end{equation}

A nonnegative balanced vector on $E$ decomposes into the twelve
circuits above. Indeed, a minimal positive dependence uses an opposite
pair or a triangle containing zero in its interior. The eight-point
square has exactly the four equal-weight and four $(1,1,2)$ triangles
listed above. Subtract the largest nonnegative multiple of an active
minimal dependence and repeat. Each subtraction removes at least one
positive coefficient, so there are at most eight steps. Write the
result as $w=\sum_C\lambda_C b^C$, where $\lambda_C\ge0$ and $b^C$
has weights $(1,1)$, $(1,1,1)$ or $(1,1,2)$.

Allocate $\lambda_C b^C_e/z_c^e$ to circuit $C$ and set
$s_C=\lambda_C\sum_e b^C_e$. Each circuit attains equality at $z_c$,
so it meets exactly the corresponding geometric-mean bound. For an
ordinary triangle choose $x=\sqrt{r_1r_2}$ and
$y=\sqrt{r_3t}$; the equality $t^3=r_1r_2r_3$ makes all three
blocks positive semidefinite. For a weighted triangle choose
$x=\sqrt{r_1r_2}$; its second block has determinant zero.
The pair block also has determinant zero. Zero circuit weight gives
zero blocks without any division by that weight.

The off-face coefficients have not been used. Assign each to its
opposite-pair allocation; this only increases the corresponding
diagonal entry. Thus the pair allocations are exactly
\eqref{eq:residual}. Equation~\eqref{eq:balance} gives
\eqref{eq:constant}, proving the converse even when the original
Laurent polynomial has no finite minimizer.

There are 24 triangle allocations, 12 constants $s_C$, eight
variables for the four ordinary triangles and four for the weighted
triangles: 48 auxiliaries in total. The pair, ordinary-triangle and
weighted-triangle blocks number $4+12+8=24$.
\end{proof}

The accompanying \texttt{ANGULAR\_CONE\_ORDER49.coo.json} lists all
109 nonzero rational coefficient records of the pencil. Off-diagonal
records are mirrored without a factor of one half. A standard-library
checker reconstructs every coefficient matrix and checks exact boundary,
zero-weight and off-face witnesses.

\section{Application to canonical Bell inequalities}

For nonnegative canonical coefficients
$c=(d,\lambda,n,m,C,A,D,\alpha,\beta)$, the state determinant is
\[
 \det P_{\rm aff}(u,v)
 =16(1+u^2)^2(1+v^2)^2u^2v^2
       \left(h(c)+\sum_{e\in E}q_e(c)u^{2e_1}v^{2e_2}\right).
\]
The nine polynomials $h(c),q_e(c)$ are given in the accompanying
determinant certificate and coefficient file. All outer coefficients
are nonnegative on this orthant. If $d>0$ and $uv\ne0$, the
three-dimensional transverse state block is positive definite.
The Schur complement therefore shows that $P_{\rm aff}\pos$ exactly
when its determinant is nonnegative. Continuity includes the testing
axes and projective endpoints. When $d=0$, the full state polynomial
is already a nonnegative sum of event squares.

Consequently a canonical Bell inequality is quantum-valid if and only
if the order-49 pencil is feasible at $(h(c),q(c))$. This gives a smaller
exact validity test than completing the entire state polynomial.
The coefficient map is quartic in $c$. The pencil is affine in the
nine Laurent coefficients; it is not an affine lift in the Bell
coefficients or the eight quantum expectations. The separate order-133
construction supplies the fixed square factors for the determinant-
weighted full state polynomial.
\end{document}
